\usemodule[zhfonts]
\input snail
\startuseMPgraphic{foo}
尺寸(里, 1cm);
宫(甲, "初始状态");
宫(乙, "可能是续行符"); 定位(乙, 乙.酉门 位于 甲.卯门 偏 (东 2里)); % 基于门的相对位置进行定位
宫(丙, "代码片段") 位于 乙 偏 (东 5里 + 北 2里);
宫(丁, "文档片段") 位于 乙 偏 (东 5里 + 南 1.5里);
呜呼 甲, 乙, 丙, 丁;

forsuffixes i = 甲, 乙, 丙, 丁:
    中央四角十二门(i, 显); 郊(i, 显);
endfor;

路(甲甲, 甲.子门 纬至 甲.坤 经至 甲.艮 经转 甲.酉门) 标注("其他字符", "北");
路(甲丙, 甲.丑门 经转 丙.酉门) 标注("其他字符", "北");
路(甲乙, 甲.卯门 => 乙.酉门) 标注("{\tt\backslash}", "北");
路(乙乙, 乙.子门 纬至 乙.坤 经至 乙.震 经转 乙.寅门) 标注("空格或制表符", "北");
路(乙甲, 乙.未门 向 (南行 纬距(乙.午门, 丁.酉门)) 纬转 甲.午门) 标注("{\tt\backslash n}", "北");
路(乙丙, 乙.卯门 纬转 丙.午门) 标注("{\tt\char`\#}", "北");
路(乙丁, 乙.午门 经转 丁.酉门) 标注("其他字符", "北");
\stopuseMPgraphic

\startTEXpage[offset=4pt]
\useMPgraphic{foo}
\stopTEXpage
